\documentclass[11pt]{article}

\usepackage{amsmath}
\usepackage{amssymb}
\usepackage{amsthm}
\usepackage{physics}
\usepackage[margin=1in]{geometry}
\usepackage{enumitem}

\newcommand{\Cinf}{C^\infty(\mathcal M)}
\newcommand{\TpM}{T_p\mathcal M}
\newcommand{\TpsM}{T_p^*\mathcal M}

\title{Mock Exam --- Fields on a Surface in $\mathbb R^3$\\ \large Based on the worked example of Appendix C.4}
\author{From Differential Geometry to General Relativity --- Study Companion}
\date{}

% ---------------------------------------------------------------------------
% Problem statements are stored as macros so that the solutions can restate
% them verbatim.
% ---------------------------------------------------------------------------

\newcommand{\probOne}{%
\begin{enumerate}[label=(\alph*)]
\item \textbf{[5 pts]} Let $\mathcal W:=\mathcal M\setminus\{(0,0,h(0,0))\}$, and let $\widehat{\mathcal Y}:\mathcal W\to\mathbb R^2$ be
defined by the same formula as $\mathcal Y$. Show that $\mathcal Y(\mathcal V)$ is open in $\mathbb R^2$, and that $\widehat{\mathcal Y}$
is not continuous at the point $r:=(-2,0,h(-2,0))$. Conclude why the chart is defined on $\mathcal V$ and not on $\mathcal W$.
\item \textbf{[5 pts]} Write down the transition maps $\mathcal T_{\mathcal{UV}}$ and $\mathcal T_{\mathcal{VU}}$ together with their
domains, and prove that both are $C^\infty$. Why does this make the atlas $C^\infty$-compatible?
\item \textbf{[5 pts]} Let $q:=(-1,-1,c_q)\in\mathcal M$. Find $c_q$, $\mathcal X(q)$ and $\mathcal Y(q)$. Then find all three
cartesian coordinates of the point $s\in\mathcal V$ with $\mathcal Y(s)=(2,-\pi/2)$.
\end{enumerate}}

\newcommand{\probTwo}{%
Let $p\in\mathcal M$ be the point with $\mathcal X(p)=(1,1)$.
\begin{enumerate}[label=(\alph*)]
\item \textbf{[4 pts]} Compute $\mathcal Y(p)$, the matrices $J_{\mathcal{VU}}(\mathcal X(p))$ and $J_{\mathcal{UV}}(\mathcal Y(p))$, and verify that
they are inverse to each other.
\item \textbf{[4 pts]} Express $(\partial x_1)_p$ and $(\partial x_2)_p$ in terms of $(\partial y_1)_p$ and $(\partial y_2)_p$.
\item \textbf{[4 pts]} Let $V_p:=2(\partial x_1)_p-(\partial x_2)_p\in T_p\mathcal M$. Find its components in the chart $(\mathcal V,\mathcal Y)$, and
verify the result by substituting the answer of (b).
\item \textbf{[4 pts]} Let $\omega_p:=3(dx^1)_p+(dx^2)_p\in T_p^*\mathcal M$. Find its components in $(\mathcal V,\mathcal Y)$, and check that
$\omega_p(V_p)$ is the same number when computed in either chart.
\item \textbf{[4 pts]} Prove that this is not a coincidence: for arbitrary $\omega_p\in T_p^*\mathcal M$ and $V_p\in T_p\mathcal M$, the
contraction $\omega_p(V_p)$ computed from $\mathcal X$-components equals the one computed from $\mathcal Y$-components.
\end{enumerate}}

\newcommand{\probThree}{%
Let $g:\mathcal M\to\mathbb R$ be the smooth scalar field
\[
g(a,b,c):=a^2+b^2.
\]
Let $p$ be the point with $\mathcal X(p)=(1,1)$ (as in Problem 2) and consider the two curves
\begin{align*}
\sigma:\mathbb R\to\mathcal M,\qquad \sigma(\lambda)&:=(\lambda,\lambda,h(\lambda,\lambda)),\\
\delta:\mathbb R\to\mathcal M,\qquad \delta(\lambda)&:=(\sqrt2\cos\lambda,\ \sqrt2\sin\lambda,\ h(\sqrt2\cos\lambda,\sqrt2\sin\lambda)).
\end{align*}
Both pass through $p$: $\sigma$ at $\lambda=1$ and $\delta$ at $\lambda=\pi/4$. Write $\sigma_p:=v_{\sigma,p}$ and $\delta_p:=v_{\delta,p}$ for their velocities at $p$.
\begin{enumerate}[label=(\alph*)]
\item \textbf{[5 pts]} Compute the components of $\sigma_p$ and $\delta_p$ in both charts, directly from the chart representatives of the
curves, and check them against the transformation law of vector components.
\item \textbf{[5 pts]} Determine for which $\lambda$ the curve $\sigma$ lies in $\mathcal V$, and describe $\mathcal Y\circ\sigma$. Is $\sigma$ smooth
as a curve on $\mathcal M$?
\item \textbf{[5 pts]} Compute the components of $(dg)_p$ in both charts from the chart representatives of $g$, and check them against the
transformation law of covector components.
\item \textbf{[5 pts]} Compute $\sigma_p(g)$ and $\delta_p(g)$ \emph{without any chart}, from the definition $v_{\gamma,p}(g):=(g\circ\gamma)'(\lambda_0)$,
and confirm both numbers by contracting components in each chart. Interpret the value of $\delta_p(g)$ geometrically.
\end{enumerate}}

\newcommand{\probFour}{%
Besides $g$ of Problem 3, consider the three vector fields $\mathfrak X,\mathfrak Y,\mathfrak R\in\Gamma T\mathcal M$ whose restrictions to $\mathcal U$ are
\[
\mathfrak X|_{\mathcal U}=\partial x_2,\qquad
\mathfrak Y|_{\mathcal U}=-b\,\partial x_1+a\,\partial x_2,\qquad
\mathfrak R|_{\mathcal U}=a\,\partial x_1+b\,\partial x_2,
\]
where $(a,b,c)$ denote the coordinates of the point at which the component functions are evaluated.
(The first two are the fields of vertical lines and of circles of the worked example, and the third one is new.)
\begin{enumerate}[label=(\alph*)]
\item \textbf{[5 pts]} Compute $\mathfrak R|_{\mathcal V}$, write $\mathfrak R$ in the tangent-bundle chart $\xi_{\mathcal Y}$, and say where $\mathfrak R$ vanishes.
\item \textbf{[5 pts]} Compute the functions $\mathfrak X[g]$, $\mathfrak Y[g]$ and $\mathfrak R[g]$ on $\mathcal M$, and interpret each of them.
\item \textbf{[5 pts]} Compute $(g\odot_\Gamma\mathfrak X\oplus_\Gamma\mathfrak Y)|_{\mathcal U}$, and then $dg\lceil g\odot_\Gamma\mathfrak X\oplus_\Gamma\mathfrak Y\rfloor$
in two ways: directly, and using the $C^\infty(\mathcal M)$-linearity of $dg\lceil\cdot\rfloor$.
\item \textbf{[5 pts]} Show that at every point $p\neq(0,0,h(0,0))$ the pair $(\mathfrak R_p,\mathfrak Y_p)$ is a basis of $T_p\mathcal M$, and
find the functions $c^1,c^2$ on $\mathcal M\setminus\{(0,0,h(0,0))\}$ such that $\mathfrak X=c^1\odot\mathfrak R\oplus c^2\odot\mathfrak Y$. Why can $(\mathfrak R,\mathfrak Y)$ not
be a frame on all of $\mathcal M$?
\item \textbf{[5 pts]} Compute $\mathfrak X[g\odot_{C^\infty} g]$ in two ways: directly, and via the Leibniz rule. Show that, unlike $dg\lceil\cdot\rfloor$, the map
$\mathfrak X[\cdot]$ is not $C^\infty(\mathcal M)$-linear.
\end{enumerate}}

\newcommand{\probFive}{%
Let $S:=dg\otimes dg$, i.e.\ the $(0,2)$-tensor field
\[
S:\Gamma T\mathcal M\times\Gamma T\mathcal M\to\Cinf,\qquad S(\mathfrak Z_1,\mathfrak Z_2):=dg\lceil\mathfrak Z_1\rfloor\odot_{C^\infty}dg\lceil\mathfrak Z_2\rfloor,
\]
with $g$ as in Problem 3.
\begin{enumerate}[label=(\alph*)]
\item \textbf{[5 pts]} Find the component matrices $(S_{\mathcal U})_{ij}$ and $(S_{\mathcal V})_{kl}$, and verify the transformation law of
a $(0,2)$-tensor field for the component $(S_{\mathcal V})_{11}$.
\item \textbf{[5 pts]} Compute $S(\mathfrak X,\mathfrak X)$, $S(\mathfrak X,\mathfrak R)$ and $S(\mathfrak Y,\mathfrak Z)$ for an arbitrary field $\mathfrak Z$.
\item \textbf{[5 pts]} Prove that $S$ is $C^\infty(\mathcal M)$-linear in its first argument.
\item \textbf{[5 pts]} True or false? Justify each answer in one or two sentences.
\begin{enumerate}[label=(\roman*)]
\item The number $\mathfrak Y_p(g)$ depends on which chart is used to compute it.
\item The field $\mathfrak Y$ is only defined on $\mathcal V$, because there it takes the simple form $\partial y_2$.
\item The curve $\gamma_{x_0}(\lambda)=(x_0,\lambda,h(x_0,\lambda))$ with $x_0\le 0$ is not smooth, since $\mathcal Y\circ\gamma_{x_0}$ is discontinuous.
\item Every covector field on $\mathcal M$ is the differential of a scalar field on $\mathcal M$.
\end{enumerate}
\end{enumerate}}

\begin{document}
\maketitle

\subsection*{Instructions}

\noindent Duration: 60 minutes. No notes. Total: 100 points (indicative time: P1 10 min, P2 12 min, P3 12 min, P4 15 min, P5 11 min).
Justify every answer; results obtained by reading them off from a chart without saying why receive no credit.
The exam is deliberately slightly longer than one hour of comfortable work: skip and return.

\subsection*{Setting}

Everything below takes place on the following $2$-dimensional smooth manifold $(\mathcal M,\mathcal O_{\mathcal M},\mathcal A_{C^\infty})_2$.
The set is the graph of a height function $h$:
\[
\mathcal M:=\{p=(a,b,c)\in\mathbb R^3\mid c=h(a,b)\},\qquad h(a,b):=\cos a\cos b+1,
\]
with the topology induced from $\mathbb R^3$. The atlas consists of two charts:
\begin{align*}
\mathcal U&:=\mathcal M,& \mathcal X(p)&:=(a,b),\\
\mathcal V&:=\mathcal M\setminus\{a\le 0,\ b=0\},& \mathcal Y(p)&:=\bigl(\sqrt{a^2+b^2},\ \mathrm{atan2}(b,a)\bigr),
\end{align*}
so that $\mathcal X(\mathcal U)=\mathbb R^2$ and $\mathcal Y(\mathcal V)=(0,\infty)\times(-\pi,\pi)$. Here $\mathrm{atan2}(b,a)\in(-\pi,\pi]$ is the polar angle of $(a,b)$.
Points of $\mathcal X(\mathcal U)$ are written $(\alpha_1,\alpha_2)$ and points of $\mathcal Y(\mathcal V)$ are written $(\beta_1,\beta_2)$; we abbreviate
$\rho:=\sqrt{\alpha_1^2+\alpha_2^2}$. The inverse chart maps are
\[
\mathcal X^{-1}(\alpha_1,\alpha_2)=(\alpha_1,\alpha_2,h(\alpha_1,\alpha_2)),\qquad
\mathcal Y^{-1}(\beta_1,\beta_2)=(\beta_1\cos\beta_2,\ \beta_1\sin\beta_2,\ h(\beta_1\cos\beta_2,\beta_1\sin\beta_2)).
\]

\paragraph{Facts you may use.} The Jacobians of the transition maps, with $(J_{\mathcal{VU}})^i{}_j=\partial\mathcal Y^i/\partial\mathcal X^j$
(row $i$, column $j$) and $(J_{\mathcal{UV}})^i{}_j=\partial\mathcal X^i/\partial\mathcal Y^j$, are
\[
J_{\mathcal{VU}}=\begin{bmatrix}\alpha_1/\rho & \alpha_2/\rho\\ -\alpha_2/\rho^2 & \alpha_1/\rho^2\end{bmatrix},\qquad
J_{\mathcal{UV}}=\begin{bmatrix}\cos\beta_2 & -\beta_1\sin\beta_2\\ \sin\beta_2 & \beta_1\cos\beta_2\end{bmatrix}.
\]
For $V_p\in T_p\mathcal M$ with components $x^i$ (in $\mathcal X$) and $y^i$ (in $\mathcal Y$), for $\omega_p\in T_p^*\mathcal M$ with components $\mathsf x_i$
and $\mathsf y_i$, and for the chart-induced bases $(\partial x_i)_p,(\partial y_i)_p$ of $T_p\mathcal M$ and $(dx^i)_p,(dy^i)_p$ of $T_p^*\mathcal M$
(summation over repeated indices is implied):
\begin{align*}
&(\partial x_j)_p=(J_{\mathcal{VU}})^i{}_j\,(\partial y_i)_p, &&(\partial y_j)_p=(J_{\mathcal{UV}})^i{}_j\,(\partial x_i)_p,\\
&y^i=(J_{\mathcal{VU}})^i{}_j\,x^j, &&x^i=(J_{\mathcal{UV}})^i{}_j\,y^j,\\
&\mathsf y_j=(J_{\mathcal{UV}})^i{}_j\,\mathsf x_i, &&\mathsf x_j=(J_{\mathcal{VU}})^i{}_j\,\mathsf y_i.
\end{align*}
The velocity of a curve $\gamma:\mathbb R\to\mathcal M$ at $p=\gamma(\lambda_0)$ is $v_{\gamma,p}$; its components in $\mathcal X$ are
$x^i=(\mathcal X^i\circ\gamma)'(\lambda_0)$, and it acts on $f\in\Cinf$ by
\[
v_{\gamma,p}(f):=(f\circ\gamma)'(\lambda_0).
\]
The differential $(df)_p\in T_p^*\mathcal M$ of $f\in\Cinf$ is defined by $(df)_p(v_{\gamma,p}):=v_{\gamma,p}(f)$; its components are
$\mathsf x_i=\partial_i[f\circ\mathcal X^{-1}](\mathcal X(p))$ and $\mathsf y_i=\partial_i[f\circ\mathcal Y^{-1}](\mathcal Y(p))$.
For vector fields $\mathfrak Z,\mathfrak W\in\Gamma T\mathcal M$ (smooth sections of $T\mathcal M$, written in fraktur to avoid clashing with the chart names), $f\in\Cinf$, and a covector field $\omega$:
\begin{align*}
&(\mathfrak Z\oplus_\Gamma\mathfrak W)(p):=\mathfrak Z_p+\mathfrak W_p, &&(f\odot_\Gamma\mathfrak Z)(p):=f(p)\,\mathfrak Z_p,\\
&\mathfrak Z[f](p):=\mathfrak Z_p(f), &&\omega\lceil\mathfrak Z\rfloor(p):=\omega_p(\mathfrak Z_p).
\end{align*}
Component functions are the functions of $p=(a,b,c)$ obtained by reading the components at each point, e.g.\ $\mathfrak Z|_{\mathcal U}=Z^i\,\partial x_i$.
A $(0,2)$-tensor field is a $C^\infty(\mathcal M)$-bilinear map $\Gamma T\mathcal M\times\Gamma T\mathcal M\to\Cinf$, with components
$T_{ij}:=T(\partial x_i,\partial x_j)$ in $\mathcal U$ and $T_{kl}:=T(\partial y_k,\partial y_l)$ in $\mathcal V$.

\newpage
\section*{Problem 1 --- The manifold and its charts \textnormal{(15 pts)}}
\probOne

\section*{Problem 2 --- Jacobians and transformation laws at a point \textnormal{(20 pts)}}
\probTwo

\newpage
\section*{Problem 3 --- Curves, velocities and a scalar field \textnormal{(20 pts)}}
\probThree

\section*{Problem 4 --- Vector fields and the module structure \textnormal{(25 pts)}}
\probFour

\newpage
\section*{Problem 5 --- A tensor field, and true or false \textnormal{(20 pts)}}
\probFive

% =============================================================================
\newpage
\part*{Solutions}
% =============================================================================

\section*{Solution to Problem 1}
\probOne

\paragraph{Solution.}
\begin{enumerate}[label=(\alph*)]
\item The image is a product of open intervals,
\[
\mathcal Y(\mathcal V)=(0,\infty)\times(-\pi,\pi),
\]
hence open in $\mathbb R^2$.

For the second claim, take $\varepsilon>0$ and the points
\[
r_\pm(\varepsilon):=(-2,\pm\varepsilon,h(-2,\pm\varepsilon))\in\mathcal W.
\]
Both converge to $r$ as $\varepsilon\to0^+$, and $r\in\mathcal W$ since $r\neq(0,0,h(0,0))$. However,
\[
\widehat{\mathcal Y}^2(r_\pm(\varepsilon))=\mathrm{atan2}(\pm\varepsilon,-2)=\arctan(\mp\varepsilon/2)\pm\pi\ \xrightarrow[\varepsilon\to0^+]{}\ \pm\pi,
\]
so the limits along the two sequences differ by $2\pi$: $\widehat{\mathcal Y}$ has no limit at $r$, and is discontinuous there.

Since a chart map must be a homeomorphism onto an open set, $\mathcal Y$ cannot be defined across the closed ray $\{a\le0,b=0\}$ (this is a
statement about the angle coordinate: the angle jumps by $2\pi$ when the ray is crossed). The origin's lift
$(0,0,h(0,0))=(0,0,2)$ is excluded as well, because there the radial coordinate would be $0\notin(0,\infty)$ and $\mathrm{atan2}$ is undefined.
\item Since $\mathcal U\cap\mathcal V=\mathcal V$, the maps are
\begin{align*}
\mathcal T_{\mathcal{UV}}=\mathcal X\circ\mathcal Y^{-1}&:\ \mathcal Y(\mathcal V)\to\mathcal X(\mathcal V),& (\beta_1,\beta_2)&\mapsto(\beta_1\cos\beta_2,\ \beta_1\sin\beta_2),\\
\mathcal T_{\mathcal{VU}}=\mathcal Y\circ\mathcal X^{-1}&:\ \mathcal X(\mathcal V)\to\mathcal Y(\mathcal V),& (\alpha_1,\alpha_2)&\mapsto\bigl(\rho,\ \mathrm{atan2}(\alpha_2,\alpha_1)\bigr),
\end{align*}
where
\[
\mathcal X(\mathcal V)=\mathbb R^2\setminus\{(\alpha_1,0)\mid\alpha_1\le0\}.
\]
The height $h$ drops out of both formulas, because $\mathcal X$ only keeps $(a,b)$.

$\mathcal T_{\mathcal{UV}}$ is $C^\infty$ because $\cos,\sin$ and products are. For $\mathcal T_{\mathcal{VU}}$, the first component
$\rho=\sqrt{\alpha_1^2+\alpha_2^2}$ is smooth since $\rho>0$ on $\mathcal X(\mathcal V)$. For the second, on $\mathcal X(\mathcal V)$ one has $\alpha_1+\rho>0$
(equality would force $\alpha_2=0,\alpha_1\le0$), and there
\[
\mathrm{atan2}(\alpha_2,\alpha_1)=2\arctan\frac{\alpha_2}{\alpha_1+\rho},
\]
which is a composition of smooth functions. (Check: with $\alpha_1=\rho\cos\theta$, $\alpha_2=\rho\sin\theta$, $\theta\in(-\pi,\pi)$, the argument is $\tan(\theta/2)$.)

Two charts are $C^\infty$-compatible when both transition maps are $C^\infty$ on the images of the overlap; the atlas
$\{(\mathcal U,\mathcal X),(\mathcal V,\mathcal Y)\}$ is therefore $C^\infty$-compatible, and $\mathcal A\equiv\mathcal A_{C^\infty}$.
\item For $q$:
\begin{align*}
c_q&=h(-1,-1)=\cos^2(1)+1\approx1.292,\\
\mathcal X(q)&=(-1,-1),\\
\mathcal Y(q)&=\bigl(\sqrt2,\ \mathrm{atan2}(-1,-1)\bigr)=\bigl(\sqrt2,\ -3\pi/4\bigr),
\end{align*}
since $q$ lies in the third quadrant, i.e.\ $b<0$, $a<0$, and the angle is $\arctan(1)-\pi$.

For $s$, the inverse chart gives
\begin{align*}
s&=\mathcal Y^{-1}(2,-\pi/2)=\bigl(2\cos(-\pi/2),\ 2\sin(-\pi/2),\ h(0,-2)\bigr)\\
&=(0,-2,1+\cos2)\approx(0,-2,0.584).
\end{align*}
\end{enumerate}

% -----------------------------------------------------------------------------
\section*{Solution to Problem 2}
\probTwo

\paragraph{Solution.} Here $\alpha=(1,1)$, so $\rho=\sqrt2$.
\begin{enumerate}[label=(\alph*)]
\item The polar coordinates are
\[
\mathcal Y(p)=(\beta_1,\beta_2)=(\sqrt2,\ \pi/4),
\]
and, using $\cos\beta_2=\sin\beta_2=1/\sqrt2$ and $\beta_1\sin\beta_2=\beta_1\cos\beta_2=1$,
\[
J_{\mathcal{VU}}=\begin{bmatrix}1/\sqrt2&1/\sqrt2\\-1/2&1/2\end{bmatrix},\qquad
J_{\mathcal{UV}}=\begin{bmatrix}1/\sqrt2&-1\\1/\sqrt2&1\end{bmatrix}.
\]
Their product is
\[
J_{\mathcal{UV}}J_{\mathcal{VU}}=\begin{bmatrix}\tfrac12+\tfrac12&\tfrac12-\tfrac12\\ \tfrac12-\tfrac12&\tfrac12+\tfrac12\end{bmatrix}=\begin{bmatrix}1&0\\0&1\end{bmatrix}.
\]
\item Using $(\partial x_j)_p=(J_{\mathcal{VU}})^i{}_j(\partial y_i)_p$, i.e.\ reading off the columns of $J_{\mathcal{VU}}$:
\begin{align*}
(\partial x_1)_p&=\tfrac1{\sqrt2}(\partial y_1)_p-\tfrac12(\partial y_2)_p,\\
(\partial x_2)_p&=\tfrac1{\sqrt2}(\partial y_1)_p+\tfrac12(\partial y_2)_p.
\end{align*}
\item The $\mathcal X$-components are $(x^1,x^2)=(2,-1)$, hence $y^i=(J_{\mathcal{VU}})^i{}_jx^j$ gives
\begin{align*}
y^1&=\tfrac1{\sqrt2}\cdot2+\tfrac1{\sqrt2}\cdot(-1)=\tfrac1{\sqrt2},\\
y^2&=-\tfrac12\cdot2+\tfrac12\cdot(-1)=-\tfrac32.
\end{align*}
Check with (b):
\begin{align*}
2(\partial x_1)_p-(\partial x_2)_p&=\Bigl(\tfrac2{\sqrt2}-\tfrac1{\sqrt2}\Bigr)(\partial y_1)_p+\Bigl(-1-\tfrac12\Bigr)(\partial y_2)_p\\
&=\tfrac1{\sqrt2}(\partial y_1)_p-\tfrac32(\partial y_2)_p.
\end{align*}
\item The $\mathcal X$-components are $(\mathsf x_1,\mathsf x_2)=(3,1)$, so $\mathsf y_j=(J_{\mathcal{UV}})^i{}_j\mathsf x_i$ gives
\begin{align*}
\mathsf y_1&=\tfrac1{\sqrt2}\cdot3+\tfrac1{\sqrt2}\cdot1=2\sqrt2,\\
\mathsf y_2&=-1\cdot3+1\cdot1=-2.
\end{align*}
The contractions are
\begin{align*}
\mathsf x_ix^i&=3\cdot2+1\cdot(-1)=5,\\
\mathsf y_iy^i&=2\sqrt2\cdot\tfrac1{\sqrt2}+(-2)\cdot\bigl(-\tfrac32\bigr)=2+3=5.
\end{align*}
\item Using the laws for vector and covector components,
\begin{align*}
\mathsf y_jy^j&=(J_{\mathcal{UV}})^i{}_j\,\mathsf x_i\;(J_{\mathcal{VU}})^j{}_k\,x^k\\
&=\mathsf x_i\,\bigl[(J_{\mathcal{UV}})^i{}_j(J_{\mathcal{VU}})^j{}_k\bigr]\,x^k\\
&=\mathsf x_i\,\delta^i{}_k\,x^k=\mathsf x_ix^i,
\end{align*}
because $J_{\mathcal{UV}}J_{\mathcal{VU}}=\mathbb 1$ (the two transition maps are inverse to each other, so their Jacobians are inverse matrices by
the chain rule). The covector components transform with $J_{\mathcal{UV}}$ contracted on its upper index, and the vector components with
$J_{\mathcal{VU}}$; this is exactly what makes the pairing cancel the chart dependence.
\end{enumerate}

% -----------------------------------------------------------------------------
\section*{Solution to Problem 3}
\probThree

\paragraph{Solution.}
\begin{enumerate}[label=(\alph*)]
\item For $\sigma$ at $\lambda_0=1$, the chart representatives are
\begin{align*}
(\mathcal X\circ\sigma)(\lambda)&=(\lambda,\lambda),\\
(\mathcal Y\circ\sigma)(\lambda)&=(\sqrt2\lambda,\ \pi/4)\qquad(\lambda>0),
\end{align*}
so the components of $\sigma_p$ are
\[
(x^1,x^2)=(1,1),\qquad(y^1,y^2)=(\sqrt2,0).
\]
Transformation law:
\[
J_{\mathcal{VU}}\begin{pmatrix}1\\1\end{pmatrix}=\begin{pmatrix}\tfrac1{\sqrt2}+\tfrac1{\sqrt2}\\-\tfrac12+\tfrac12\end{pmatrix}=\begin{pmatrix}\sqrt2\\0\end{pmatrix}.\ \checkmark
\]
For $\delta$ at $\lambda_0=\pi/4$, near $\lambda_0$ the curve stays in $\mathcal V$, and
\begin{align*}
(\mathcal Y\circ\delta)(\lambda)&=(\sqrt2,\ \lambda),\\
(\mathcal X\circ\delta)(\lambda)&=(\sqrt2\cos\lambda,\ \sqrt2\sin\lambda),
\end{align*}
so the components of $\delta_p$ are
\[
(y^1,y^2)=(0,1),\qquad(x^1,x^2)=(-\sqrt2\sin\tfrac\pi4,\ \sqrt2\cos\tfrac\pi4)=(-1,1).
\]
Transformation law:
\[
J_{\mathcal{UV}}\begin{pmatrix}0\\1\end{pmatrix}=\begin{pmatrix}-1\\1\end{pmatrix}.\ \checkmark
\]
Thus $\sigma_p=(\partial x_1)_p+(\partial x_2)_p=\sqrt2(\partial y_1)_p$ (purely radial) and $\delta_p=-(\partial x_1)_p+(\partial x_2)_p=(\partial y_2)_p$ (purely angular).
\item The curve leaves $\mathcal V$ exactly when $\mathcal X(\sigma(\lambda))=(\lambda,\lambda)$ satisfies $b=0$ and $a\le0$, i.e.\ only for $\lambda=0$ (the origin's lift).
So $\sigma\in\mathcal V$ for $\lambda\neq0$, and
\[
(\mathcal Y\circ\sigma)(\lambda)=\begin{cases}(\sqrt2\lambda,\ \pi/4),&\lambda>0,\\ (\sqrt2|\lambda|,\ -3\pi/4),&\lambda<0,\end{cases}
\]
two disconnected half-lines in $\mathcal Y(\mathcal V)$. Nevertheless $\sigma$ is smooth: smoothness is tested in a chart that sees the whole curve, and
$\mathcal X\circ\sigma=(\lambda,\lambda)$ is smooth on $\mathbb R$. The gap belongs to the chart $\mathcal Y$, not to the curve.
\item The chart representatives of $g$ are
\begin{align*}
g\circ\mathcal X^{-1}&:(\alpha_1,\alpha_2)\mapsto\alpha_1^2+\alpha_2^2,\\
g\circ\mathcal Y^{-1}&:(\beta_1,\beta_2)\mapsto\beta_1^2\cos^2\beta_2+\beta_1^2\sin^2\beta_2=\beta_1^2.
\end{align*}
Hence, at $\alpha=(1,1)$ and $\beta=(\sqrt2,\pi/4)$,
\begin{align*}
(\mathsf x_1,\mathsf x_2)&=(2\alpha_1,2\alpha_2)=(2,2),\\
(\mathsf y_1,\mathsf y_2)&=(2\beta_1,0)=(2\sqrt2,0).
\end{align*}
Transformation law $\mathsf y_j=(J_{\mathcal{UV}})^i{}_j\mathsf x_i$:
\begin{align*}
\mathsf y_1&=\tfrac1{\sqrt2}\cdot2+\tfrac1{\sqrt2}\cdot2=2\sqrt2,\\
\mathsf y_2&=-1\cdot2+1\cdot2=0.\ \checkmark
\end{align*}
\item Along the curves,
\begin{align*}
(g\circ\sigma)(\lambda)&=2\lambda^2,\\
(g\circ\delta)(\lambda)&=2\cos^2\lambda+2\sin^2\lambda=2,
\end{align*}
so
\begin{align*}
\sigma_p(g)&=(g\circ\sigma)'(1)=4,\\
\delta_p(g)&=(g\circ\delta)'(\pi/4)=0.
\end{align*}
Contractions in $\mathcal X$:
\begin{align*}
x^i\mathsf x_i\big|_{\sigma_p}&=1\cdot2+1\cdot2=4,\\
x^i\mathsf x_i\big|_{\delta_p}&=(-1)\cdot2+1\cdot2=0.
\end{align*}
Contractions in $\mathcal Y$:
\begin{align*}
y^i\mathsf y_i\big|_{\sigma_p}&=\sqrt2\cdot2\sqrt2+0\cdot0=4,\\
y^i\mathsf y_i\big|_{\delta_p}&=0\cdot2\sqrt2+1\cdot0=0.
\end{align*}
The four numbers agree with the chart-free ones. Since $g$ is constant along $\delta$ (the circle $a^2+b^2=2$ is a level set of $g$), the velocity
$\delta_p$ is tangent to a level set of $g$, hence $(dg)_p(\delta_p)=0$; equivalently $\delta_p\in\ker(dg)_p$. The curve $\sigma$ is radial and climbs $g$.
\end{enumerate}

% -----------------------------------------------------------------------------
\section*{Solution to Problem 4}
\probFour

\paragraph{Solution.}
\begin{enumerate}[label=(\alph*)]
\item The $\mathcal X$-component functions of $\mathfrak R$ are $(R^1,R^2)=(a,b)$, so the transformation law $y^i=(J_{\mathcal{VU}})^i{}_jx^j$ with
$\alpha=(a,b)$, $\rho=\sqrt{a^2+b^2}$ gives
\begin{align*}
\tilde Y^1&=\frac a\rho\,a+\frac b\rho\,b=\rho,\\
\tilde Y^2&=-\frac b{\rho^2}\,a+\frac a{\rho^2}\,b=0,
\end{align*}
that is,
\[
\mathfrak R|_{\mathcal V}=\sqrt{a^2+b^2}\;\partial y_1.
\]
In the tangent-bundle chart $\xi_{\mathcal Y}$ the point $\mathfrak R_p$ is described by its base point and its components,
\[
\xi_{\mathcal Y}(\mathfrak R_p)=\begin{pmatrix}(\sqrt{a^2+b^2},\ \mathrm{atan2}(b,a))\\ (\sqrt{a^2+b^2},\ 0)\end{pmatrix}.
\]
The field is smooth on all of $\mathcal M$ (its $\mathcal X$-components are smooth on $\mathcal U=\mathcal M$), and $\mathfrak R_p=0$ if and only if $a=b=0$, i.e.\ only at the lift of the origin.
\item With the components $\mathsf x_i=(2a,2b)$ of $dg$ in $\mathcal U$:
\begin{align*}
\mathfrak X[g]&=X^i\mathsf x_i=0\cdot2a+1\cdot2b=2b,\\
\mathfrak Y[g]&=\tilde X^i\mathsf x_i=(-b)(2a)+a(2b)=0,\\
\mathfrak R[g]&=a\cdot2a+b\cdot2b=2(a^2+b^2)=2g.
\end{align*}
(In $\mathcal V$: $\mathfrak Y[g]=0\cdot2\rho+1\cdot0=0$ and $\mathfrak R[g]=\rho\cdot2\rho+0\cdot0=2\rho^2$, the same functions.)
Interpretation: $\mathfrak Y$ is tangent to the circles $g=\mathrm{const}$, so it never changes $g$. $\mathfrak R$ points outward and $\mathfrak R[g]=2g$ says that $g$ grows
quadratically along it (an Euler-type relation for a homogeneous function of degree $2$). $\mathfrak X$ climbs $g$ where $b>0$, descends where $b<0$, and runs along a level set exactly where $b=0$.
\item In $\mathcal U$,
\begin{align*}
(g\odot_\Gamma\mathfrak X\oplus_\Gamma\mathfrak Y)|_{\mathcal U}&=g\,\partial x_2-b\,\partial x_1+a\,\partial x_2\\
&=-b\,\partial x_1+(a^2+b^2+a)\,\partial x_2.
\end{align*}
Directly,
\begin{align*}
dg\lceil g\odot_\Gamma\mathfrak X\oplus_\Gamma\mathfrak Y\rfloor&=2a(-b)+2b(a^2+b^2+a)\\
&=2b(a^2+b^2)=2b\,g.
\end{align*}
Via $C^\infty(\mathcal M)$-linearity,
\begin{align*}
dg\lceil g\odot_\Gamma\mathfrak X\oplus_\Gamma\mathfrak Y\rfloor&=g\,dg\lceil\mathfrak X\rfloor+dg\lceil\mathfrak Y\rfloor\\
&=g\cdot2b+0=2b\,g.
\end{align*}
The two agree, as they must: at each $p$, $(dg)_p$ is a linear map on $T_p\mathcal M$ and $g(p)$ is just a real number multiplying the vector $\mathfrak X_p$.
\item In the $\mathcal X$-components the matrix whose columns are $\mathfrak R_p$ and $\mathfrak Y_p$ is
\[
M=\begin{bmatrix}a&-b\\b&a\end{bmatrix},\qquad\det M=a^2+b^2,
\]
which is nonzero unless $a=b=0$. Thus $(\mathfrak R_p,\mathfrak Y_p)$ is a basis of $T_p\mathcal M$ for every $p$ other than the lift of the origin.
The coefficients solve $M\binom{c^1}{c^2}=\binom01$, where $(0,1)$ are the $\mathcal X$-components of $\mathfrak X$:
\[
\binom{c^1}{c^2}=M^{-1}\binom01=\frac1{a^2+b^2}\begin{bmatrix}a&b\\-b&a\end{bmatrix}\binom01=\binom{b/(a^2+b^2)}{a/(a^2+b^2)}.
\]
Hence, on $\mathcal M\setminus\{(0,0,h(0,0))\}$,
\[
\mathfrak X=\frac{b}{a^2+b^2}\odot_\Gamma\mathfrak R\ \oplus_\Gamma\ \frac a{a^2+b^2}\odot_\Gamma\mathfrak Y.
\]
(Check in $\mathcal V$: $\frac b{\rho^2}\cdot\rho\,\partial y_1+\frac a{\rho^2}\,\partial y_2=\frac b\rho\,\partial y_1+\frac a{\rho^2}\,\partial y_2$, the known form of $\mathfrak X|_{\mathcal V}$.)

At the lift of the origin both $\mathfrak R$ and $\mathfrak Y$ vanish, so they span nothing in the $2$-dimensional space $T_p\mathcal M$; moreover $c^1,c^2$ blow up there and are not
elements of $C^\infty(\mathcal M)$. The coefficients are non-constant functions: linear combinations in $\Gamma T\mathcal M$ have scalars in $C^\infty(\mathcal M)$, not in $\mathbb R$, which is the
module structure at work; and a frame that exists over an open subset need not extend to all of $\mathcal M$.
\item Along $\gamma_a(\lambda)=(a,\lambda,h(a,\lambda))$, at $\lambda_0=b$, directly,
\[
\mathfrak X[g\odot g]=\frac{d}{d\lambda}\Bigl[(a^2+\lambda^2)^2\Bigr]_{\lambda=b}=4b(a^2+b^2).
\]
Via the Leibniz rule and (b),
\[
\mathfrak X[g\odot g]=g\odot\mathfrak X[g]\oplus g\odot\mathfrak X[g]=2g\cdot2b=4b(a^2+b^2).\ \checkmark
\]
If $\mathfrak X[\cdot]$ were $C^\infty(\mathcal M)$-linear we would have $\mathfrak X[g\odot g]=g\odot\mathfrak X[g]$. The difference of the two sides is
\[
\mathfrak X[g\odot g]-g\odot\mathfrak X[g]=g\,\mathfrak X[g]=2b(a^2+b^2),
\]
nonzero for instance at $a=b=1$ (where $\mathfrak X[g\odot g]=8$ but $g\,\mathfrak X[g]=4$). So $\mathfrak X[\cdot]$ is only $\mathbb R$-linear plus Leibniz, whereas $dg\lceil\cdot\rfloor$ is $C^\infty(\mathcal M)$-linear
(part (c)); this is why $dg$, not $\mathfrak X[\cdot]$, defines a tensor field.
\end{enumerate}

% -----------------------------------------------------------------------------
\section*{Solution to Problem 5}
\probFive

\paragraph{Solution.}
\begin{enumerate}[label=(\alph*)]
\item By definition $(S_{\mathcal U})_{ij}=S(\partial x_i,\partial x_j)=\mathsf x_i\mathsf x_j$ with $\mathsf x_i=(2a,2b)$, and likewise $(S_{\mathcal V})_{kl}=\mathsf y_k\mathsf y_l$ with $\mathsf y_k=(2\rho,0)$
(where $\rho=\sqrt{a^2+b^2}$, as in Problem 3(c) with $\beta_1=\rho$):
\[
(S_{\mathcal U})_{ij}=\begin{bmatrix}4a^2&4ab\\4ab&4b^2\end{bmatrix},\qquad
(S_{\mathcal V})_{kl}=\begin{bmatrix}4(a^2+b^2)&0\\0&0\end{bmatrix}.
\]
The transformation law is $(S_{\mathcal V})_{kl}=(J_{\mathcal{UV}})^i{}_k(J_{\mathcal{UV}})^j{}_l(S_{\mathcal U})_{ij}$. For $k=l=1$, with $a=\beta_1\cos\beta_2$ and $b=\beta_1\sin\beta_2$,
\begin{align*}
(S_{\mathcal V})_{11}&=\cos^2\beta_2\,(4a^2)+2\cos\beta_2\sin\beta_2\,(4ab)+\sin^2\beta_2\,(4b^2)\\
&=4(a\cos\beta_2+b\sin\beta_2)^2\\
&=4\beta_1^2(\cos^2\beta_2+\sin^2\beta_2)^2=4\beta_1^2=4(a^2+b^2).\ \checkmark
\end{align*}
\item Using $S(\mathfrak Z_1,\mathfrak Z_2)=dg\lceil\mathfrak Z_1\rfloor\,dg\lceil\mathfrak Z_2\rfloor$ and $dg\lceil\mathfrak Z\rfloor=\mathfrak Z[g]$ from Problem 4(b):
\begin{align*}
S(\mathfrak X,\mathfrak X)&=(2b)^2=4b^2,\\
S(\mathfrak X,\mathfrak R)&=2b\cdot2g=4b(a^2+b^2),\\
S(\mathfrak Y,\mathfrak Z)&=\mathfrak Y[g]\,\mathfrak Z[g]=0\cdot\mathfrak Z[g]=0\quad\text{for every }\mathfrak Z\in\Gamma T\mathcal M.
\end{align*}
So $\mathfrak Y$ lies in the kernel of $S$ at every point: $S$ is degenerate (its component matrices have rank $\le1$, with zero determinant in either chart).
\item Let $\varphi\in\Cinf$ and $\mathfrak Z_1,\mathfrak W_1,\mathfrak Z_2\in\Gamma T\mathcal M$. Since $dg\lceil\cdot\rfloor$ is $C^\infty(\mathcal M)$-linear (at each $p$, $(dg)_p$ is linear on $T_p\mathcal M$ and $\varphi(p)$ is a number),
\begin{align*}
S(\varphi\odot_\Gamma\mathfrak Z_1\oplus_\Gamma\mathfrak W_1,\mathfrak Z_2)&=dg\lceil\varphi\odot_\Gamma\mathfrak Z_1\oplus_\Gamma\mathfrak W_1\rfloor\odot_{C^\infty}dg\lceil\mathfrak Z_2\rfloor\\
&=\bigl(\varphi\odot_{C^\infty}dg\lceil\mathfrak Z_1\rfloor\oplus_{C^\infty}dg\lceil\mathfrak W_1\rfloor\bigr)\odot_{C^\infty}dg\lceil\mathfrak Z_2\rfloor\\
&=\varphi\odot_{C^\infty}S(\mathfrak Z_1,\mathfrak Z_2)\oplus_{C^\infty}S(\mathfrak W_1,\mathfrak Z_2).
\end{align*}
Here $\oplus_{C^\infty}$ and $\odot_{C^\infty}$ denote the pointwise sum and product of functions in $\Cinf$. The same computation in the second slot follows from symmetry of $S$.
\item
\begin{enumerate}[label=(\roman*)]
\item \textbf{False.} The components $y^i$ and $\mathsf y_i$ (and the partial derivatives) depend on the chart, but their contraction does not (Problem 2(e)); in Problem 3(d) both charts gave the same number.
\item \textbf{False.} $\mathfrak Y$ is a smooth vector field on all of $\mathcal M$: its $\mathcal X$-components $(-b,a)$ are smooth on $\mathcal U=\mathcal M$. Only the expression $\partial y_2$ is restricted to $\mathcal V$, since $\partial y_2$ itself only exists there.
\item \textbf{False.} Smoothness of a curve is tested in a chart that contains its image; $\mathcal X\circ\gamma_{x_0}=(x_0,\lambda)$ is smooth on $\mathbb R$. The chart $\mathcal Y$ simply cannot see the points on the ray $\{a\le0,b=0\}$, so the curve appears in it as two pieces.
\item \textbf{False.} If $\omega=d\varphi$, its component functions in $\mathcal U$ are $\mathsf x_i=\partial_i[\varphi\circ\mathcal X^{-1}]$, and equality of mixed partials forces
\[
\partial_1\mathsf x_2=\partial_2\mathsf x_1.
\]
The covector field $\omega:=a\,dx^2$ has $\mathsf x_2=a$ and $\mathsf x_1=0$, so $\partial_1\mathsf x_2=1\neq0=\partial_2\mathsf x_1$, and it is not a differential. Only \emph{closed} covector fields can be differentials.
That closed implies exact holds on $\mathcal M$ because it is a graph over $\mathbb R^2$, hence contractible; on a manifold with nontrivial first de Rham cohomology it fails even for closed fields: on $\mathbb R^2\setminus\{0\}$ the field
$\bigl(-y\,dx+x\,dy\bigr)/(x^2+y^2)$ is locally the differential of the polar angle but not of any global scalar field.
\end{enumerate}
\end{enumerate}

\end{document}
